\documentclass[12pt,a4paper]{article}
\usepackage[T2A]{fontenc}
\usepackage[utf8]{inputenc}
\usepackage[russian]{babel}
\usepackage{amsmath,amssymb}
\usepackage{geometry}
\usepackage{booktabs}
\usepackage{hyperref}
\geometry{margin=2cm}

\title{Отчёт по реализации программного средства\\для задачи о назначениях с ограничениями}
\author{Проект \texttt{assignment-optimizer}}
\date{\today}

\begin{document}
\maketitle

\section{Цель работы}
Разработать программное средство для приближённого решения задачи о назначениях с дополнительными ограничениями вида
\begin{equation}
\sum_{i=1}^{n}\sum_{j=1}^{n} d_{ij}^{k}x_{ij} - b_k \le 0,\quad k=1,\dots,K,
\end{equation}
где $X=\{x_{ij}\}$ --- матрица назначений, $x_{ij}\in\{0,1\}$, и выполнены стандартные условия задачи о назначениях:
каждый исполнитель назначается ровно на одну работу и каждая работа получает ровно одного исполнителя.

\section{Математическая постановка}
Минимизируется целевая функция
\begin{equation}
F(X)=\sum_{i=1}^{n}\sum_{j=1}^{n} c_{ij}x_{ij}\rightarrow\min,
\end{equation}
при ограничениях на назначения и дополнительных ограничениях по матрицам $D_k$.

Используется функция Лагранжа:
\begin{equation}
L(X,\lambda)=\sum_{i,j} c_{ij}x_{ij}+\sum_{k=1}^{K}\lambda_k\left(\sum_{i,j} d_{ij}^{k}x_{ij}-b_k\right),\quad \lambda_k\ge 0.
\end{equation}

На каждой итерации строится модифицированная матрица стоимости:
\begin{equation}
G = C + \sum_{k=1}^{K}\lambda_k D_k,
\end{equation}
после чего решается стандартная задача о назначениях для $G$.

\section{Реализованный алгоритм (метод Удзавы)}
\begin{enumerate}
  \item Инициализация: ввод $n, K, C, D_k, b_k, \lambda^0$, параметров остановки и шага.
  \item Решение задачи о назначениях для матрицы $G$ (алгоритм Венгера в реализации \texttt{scipy.optimize.linear\_sum\_assignment}).
  \item Вычисление субградиента:
  \[
    f_k(X)=\sum_{i,j} d_{ij}^{k}x_{ij}-b_k,\; k=1,\dots,K.
  \]
  \item Проверка остановки: $\|f(X)\|\le\varepsilon$ или достигнут лимит итераций.
  \item Обновление множителей:
  \[
  \lambda^{t+1}=\left[\lambda^{t}+\psi_t f(X^t)\right]^+,
  \]
  где $[\cdot]^+$ --- проекция на неотрицательный ортант.
\end{enumerate}

Поддержаны два режима шага:
\begin{itemize}
  \item гармонический: $\psi_t=\frac{1}{t+1}$;
  \item постоянный: $\psi_t=\text{const}$.
\end{itemize}

\section{Архитектура программного средства}
Решение выполнено в виде клиент--серверного приложения.

\subsection{Backend (Python + Flask)}
\begin{itemize}
  \item Модуль \texttt{backend/solver.py}:
  \begin{itemize}
    \item валидация входных данных;
    \item итеративная процедура Удзавы;
    \item выбор лучшего допустимого решения (или наилучшего по норме субградиента как fallback);
    \item формирование истории итераций.
  \end{itemize}
  \item HTTP API:
  \begin{itemize}
    \item \texttt{POST /api/solve} --- запуск расчёта;
    \item возврат JSON с назначением, значением целевой функции, субградиентом и историей.
  \end{itemize}
\end{itemize}

\subsection{Frontend (React + Vite)}
Пользовательский интерфейс предоставляет:
\begin{itemize}
  \item ввод размерности $n$ и числа ограничений $K$;
  \item редактирование матрицы $C$;
  \item редактирование набора матриц $D_k$;
  \item ввод векторов $b$ и $\lambda^0$;
  \item настройку параметров \texttt{max\_iter}, \texttt{eps}, режима шага;
  \item запуск вычислений и отображение результатов.
\end{itemize}

\section{Формат входных данных API}
\begin{verbatim}
{
  "n": 4,
  "k": 2,
  "c": [[...], ...],
  "d": [[[...]], [[...]]],
  "b": [b1, b2],
  "lambda0": [0, 0],
  "max_iter": 200,
  "eps": 1e-6,
  "step_mode": "harmonic",
  "step_const": 0.1
}
\end{verbatim}

\section{Формат выходных данных API}
\begin{verbatim}
{
  "best_assignment_zero_based": [...],
  "best_assignment_one_based": [...],
  "best_objective": ...,
  "best_subgradient": [...],
  "best_iteration": ...,
  "iterations_done": ...,
  "history": [...]
}
\end{verbatim}

\section{Проверка корректности и ограничения}
\begin{itemize}
  \item Проверяются размеры: $C\in\mathbb{R}^{n\times n}$, $D\in\mathbb{R}^{K\times n\times n}$, $b\in\mathbb{R}^{K}$, $\lambda^0\in\mathbb{R}^{K}$.
  \item Гарантируется проекция множителей Лагранжа на область $\lambda\ge 0$.
  \item Из-за субградиентного характера метод может сходиться медленно; качество зависит от выбора шага и числа итераций.
\end{itemize}

\section{Инструкция по запуску}
\subsection*{Backend}
\begin{verbatim}
python -m venv .venv
source .venv/bin/activate
pip install -r requirements.txt
python server.py
\end{verbatim}

\subsection*{Frontend}
\begin{verbatim}
npm install
npm run dev
\end{verbatim}

\section{Вывод}
В рамках работы реализовано программное средство для численного решения задачи о назначениях с дополнительными ограничениями на основе двойственного субградиентного подхода (метод Удзавы). Решение включает математически обоснованный backend-алгоритм и удобный интерфейс для ввода данных и анализа результатов.

\end{document}
